\def\tituloconcepto{Puertas universales y lógica reversible}
% Si el capitulo se compila solo (sin main.tex), esta cabecera abre el documento.
% Antes de % Si el capitulo se compila solo (sin main.tex), esta cabecera abre el documento.
% Antes de % Si el capitulo se compila solo (sin main.tex), esta cabecera abre el documento.
% Antes de \input{cabecera} cada capitulo define \tituloconcepto.
\ifdefined\enmain\else
  \documentclass[11pt,a4paper]{article}
  \input{preambulo}
  \begin{document}
  \title{\tituloconcepto}
  \author{\autor}
  \date{\today}
  \maketitle
\fi
 cada capitulo define \tituloconcepto.
\ifdefined\enmain\else
  \documentclass[11pt,a4paper]{article}
  % Preambulo compartido por todos los documentos. Compilar con XeLaTeX/LuaLaTeX
% (o Tectonic). Si se usa Overleaf, elegir el compilador XeLaTeX.
\usepackage[spanish,es-noshorthands]{babel}
\usepackage{amsmath,amssymb}
\usepackage{booktabs,array}
\usepackage[a4paper,margin=2.5cm]{geometry}
\usepackage{xcolor}
\usepackage{listings}
\usepackage[most]{tcolorbox}
\usepackage{quantikz}
\usepackage[hidelinks]{hyperref}

% ---- Autor (editar a gusto) -------------------------------------------------
\newcommand{\autor}{Genesis Silesky}

% ---- Notacion de Dirac y atajos ---------------------------------------------
% quantikz ya define \ket y \bra; DeclareRobustCommand los redefine sin error.
\DeclareRobustCommand{\ket}[1]{\left|#1\right\rangle}
\DeclareRobustCommand{\bra}[1]{\left\langle#1\right|}
\DeclareRobustCommand{\braket}[2]{\left\langle#1\middle|#2\right\rangle}
\newcommand{\C}{\mathbb{C}}
\newcommand{\R}{\mathbb{R}}
\DeclareMathOperator{\Tr}{Tr}

% ---- Cajas de apoyo --------------------------------------------------------
% "concepto": definiciones previas que hacen falta para entender el capitulo.
\newtcolorbox{concepto}[1]{breakable,colback=blue!4,colframe=blue!55!black,
  fonttitle=\bfseries,title={Concepto previo: #1}}
% "nota": aclaraciones y advertencias.
\newtcolorbox{nota}[1]{breakable,colback=orange!6,colframe=orange!70!black,
  fonttitle=\bfseries,title={#1}}

% ---- Estilos de listados ----------------------------------------------------
\lstdefinestyle{py}{language=Python,basicstyle=\ttfamily\footnotesize,
  breaklines=true,postbreak=\mbox{\textcolor{gray}{$\hookrightarrow$}\space},
  numbers=left,numberstyle=\tiny\color{gray},frame=single,
  commentstyle=\color{green!40!black}\itshape,
  keywordstyle=\color{blue!70!black}\bfseries,
  stringstyle=\color{red!60!black},showstringspaces=false,tabsize=4,
  columns=fullflexible,keepspaces=true}
\lstdefinestyle{salida}{basicstyle=\ttfamily\footnotesize,breaklines=true,
  frame=leftline,backgroundcolor=\color{gray!8},columns=fullflexible,
  keepspaces=true}

  \begin{document}
  \title{\tituloconcepto}
  \author{\autor}
  \date{\today}
  \maketitle
\fi
 cada capitulo define \tituloconcepto.
\ifdefined\enmain\else
  \documentclass[11pt,a4paper]{article}
  % Preambulo compartido por todos los documentos. Compilar con XeLaTeX/LuaLaTeX
% (o Tectonic). Si se usa Overleaf, elegir el compilador XeLaTeX.
\usepackage[spanish,es-noshorthands]{babel}
\usepackage{amsmath,amssymb}
\usepackage{booktabs,array}
\usepackage[a4paper,margin=2.5cm]{geometry}
\usepackage{xcolor}
\usepackage{listings}
\usepackage[most]{tcolorbox}
\usepackage{quantikz}
\usepackage[hidelinks]{hyperref}

% ---- Autor (editar a gusto) -------------------------------------------------
\newcommand{\autor}{Genesis Silesky}

% ---- Notacion de Dirac y atajos ---------------------------------------------
% quantikz ya define \ket y \bra; DeclareRobustCommand los redefine sin error.
\DeclareRobustCommand{\ket}[1]{\left|#1\right\rangle}
\DeclareRobustCommand{\bra}[1]{\left\langle#1\right|}
\DeclareRobustCommand{\braket}[2]{\left\langle#1\middle|#2\right\rangle}
\newcommand{\C}{\mathbb{C}}
\newcommand{\R}{\mathbb{R}}
\DeclareMathOperator{\Tr}{Tr}

% ---- Cajas de apoyo --------------------------------------------------------
% "concepto": definiciones previas que hacen falta para entender el capitulo.
\newtcolorbox{concepto}[1]{breakable,colback=blue!4,colframe=blue!55!black,
  fonttitle=\bfseries,title={Concepto previo: #1}}
% "nota": aclaraciones y advertencias.
\newtcolorbox{nota}[1]{breakable,colback=orange!6,colframe=orange!70!black,
  fonttitle=\bfseries,title={#1}}

% ---- Estilos de listados ----------------------------------------------------
\lstdefinestyle{py}{language=Python,basicstyle=\ttfamily\footnotesize,
  breaklines=true,postbreak=\mbox{\textcolor{gray}{$\hookrightarrow$}\space},
  numbers=left,numberstyle=\tiny\color{gray},frame=single,
  commentstyle=\color{green!40!black}\itshape,
  keywordstyle=\color{blue!70!black}\bfseries,
  stringstyle=\color{red!60!black},showstringspaces=false,tabsize=4,
  columns=fullflexible,keepspaces=true}
\lstdefinestyle{salida}{basicstyle=\ttfamily\footnotesize,breaklines=true,
  frame=leftline,backgroundcolor=\color{gray!8},columns=fullflexible,
  keepspaces=true}

  \begin{document}
  \title{\tituloconcepto}
  \author{\autor}
  \date{\today}
  \maketitle
\fi


\section{Puertas universales y lógica clásica reversible}

\subsection{Objetivo}
Ver que toda puerta de un qubit se arma con rotaciones alrededor de $z$ y de $y$,
cómo se construye su versión controlada y cómo la puerta Toffoli permite
hacer cálculos de lógica clásica (AND, OR, XOR\ldots) con un computador cuántico.

\subsection{Conceptos previos}
\begin{concepto}{rotaciones y puertas controladas}
Ver capítulo~12 (rotaciones $R_y$ y $R_z$) y la sección~\ref{sec:prerreq}
(CNOT). Una puerta \emph{controlada} aplica una operación al qubit objetivo
solo si el qubit de control vale $\ket1$.
\end{concepto}

\subsection{Fundamento teórico}
Un algoritmo cuántico es un circuito: puertas conectadas en un orden
determinado. Este capítulo muestra que con puertas de uno, dos y tres qubits se
puede armar cualquier operación.

\paragraph{Puerta universal de un qubit.} Toda puerta de un qubit es una matriz
unitaria $2\times2$ y se puede escribir como una rotación alrededor de $z$,
otra alrededor de $y$ y otra alrededor de $z$, más una fase global:
\begin{equation}
  U=e^{i\alpha}R_z(\beta)\,R_y(\gamma)\,R_z(\delta),
\end{equation}
con $R_z(\theta)=e^{-i\theta Z/2}$ y $R_y(\theta)=e^{-i\theta Y/2}$. Los ángulos
$\beta,\gamma,\delta$ son los de las rotaciones sobre la esfera de Bloch.

\paragraph{Forma $U=e^{i\alpha}AXBXC$.} Se puede escribir $U$ con tres matrices
unitarias $A,B,C$ que cumplen $ABC=I$:
\begin{equation}
  A=R_z(\beta)R_y(\tfrac\gamma2),\quad
  B=R_y(-\tfrac\gamma2)R_z(-\tfrac{\delta+\beta}2),\quad
  C=R_z(\tfrac{\delta-\beta}2).
\end{equation}
Como $XR_y(\theta)X=R_y(-\theta)$ y $XR_z(\theta)X=R_z(-\theta)$, se obtiene
$e^{i\alpha}AXBXC=U$. Esto sirve para la puerta controlada: si el control vale
$\ket0$ las $X$ no actúan y se aplica $ABC=I$ (no pasa nada); si vale
$\ket1$ se aplica $AXBXC$, que es $U$ (salvo la fase $e^{i\alpha}$, que se
agrega con una puerta de fase en el control). Así, una puerta controlada
se arma con puertas de uno y dos qubits. Con $n$ controles se obtiene la puerta
$C^n(U)$, que aplica $U$ solo si todos valen $\ket1$.

\paragraph{Toffoli.} La puerta Toffoli es la $C^2(X)$: tiene dos qubits de control
y uno objetivo, y lo invierte solo cuando ambos controles valen $\ket1$.

\paragraph{Lógica clásica con puertas cuánticas.} Toda puerta lógica clásica se
puede convertir en una puerta cuántica reversible hecha con $X$ (NOT), CNOT
y Toffoli. Con las entradas $A$, $B$ y la salida $C$:
\begin{itemize}
  \item \textbf{AND:} Toffoli con $A,B$ de controles y $C$ (en $\ket0$) como
    objetivo: $C$ pasa a $\ket1$ solo si $A=B=1$.
  \item \textbf{OR:} el objetivo $C$ empieza en $\ket1$; se aplica $X$ a $A$
    y $B$, luego la Toffoli, y de nuevo $X$ a $A$ y $B$. El objetivo pasa a
    $\ket0$ solo si $A=B=0$.
  \item \textbf{NAND, NOR, XOR\ldots:} se arman con las mismas piezas (en este
    ejemplo, NAND y NOR añaden una $X$ a la salida; XOR son dos CNOT hacia $C$).
\end{itemize}

\subsection{Código (completamente comentado)}
\lstinputlisting[style=py]{../../fuentes/14_logica_reversible.py}

\subsection{Explicación celda por celda}
\begin{description}
  \item[Celda 2] Toma una unitaria al azar, calcula sus ángulos $\beta,\gamma,\delta$
    con \texttt{OneQubitEulerDecomposer} y comprueba que se reconstruye con
    $R_zR_yR_z$.
  \item[Celda 3] Construye $A,B,C$ con esos ángulos y verifica $ABC=I$ y
    $e^{i\alpha}AXBXC=U$.
  \item[Celda 4] Imprime la tabla de verdad de la Toffoli.
  \item[Celda 5] Define el circuito de cada puerta lógica.
  \item[Celda 6] Compara cada circuito con la tabla de verdad clásica, para las
    cuatro combinaciones de $A$ y $B$.
\end{description}

\subsection{Salida obtenida}
\lstinputlisting[style=salida]{../../fuentes/salidas/14_logica_reversible.txt}
Es determinista. Los ángulos dependen de la unitaria aleatoria elegida, que usa
semilla fija.

\subsection{Interpretación}
Las dos descomposiciones se verifican (\texttt{True}). En la Toffoli, la cadena
se lee como $q_2q_1q_0$: solo cambian $011\to111$ y $111\to011$, es decir, el
objetivo $q_2$ se invierte cuando $q_0=q_1=1$. Las seis puertas lógicas
coinciden con su tabla de verdad clásica.

\subsection{Ejercicios}
\begin{enumerate}
  \item Cambia la unitaria aleatoria (otra \texttt{seed}) y verifica de nuevo las
    dos descomposiciones.
  \item Construye una puerta lógica con tres entradas, por ejemplo $C=A\wedge B\wedge D$
    (necesita una puerta con tres controles).
  \item ¿Por qué el circuito OR necesita que el objetivo empiece en $\ket1$?
\end{enumerate}

\subsection{Errores comunes}
\begin{itemize}
  \item Olvidar el orden de bits de Qiskit: la cadena \texttt{'011'} es
    $q_2=0,\,q_1=1,\,q_0=1$.
  \item Pensar que la puerta lógica clásica se aplica con una sola entrada: en
    la versión reversible las entradas se conservan (en $A$ y $B$) y la salida
    se escribe en un qubit aparte.
\end{itemize}

\ifdefined\enmain\else
  \end{document}
\fi

